\clearpage
\fancyhead[L]{\small PHILOSOPHICAL TRANSACTIONS.}
\newcommand{\ObservedWollastonDenominator}{\hbox{\vrule width0pt height4.600000bp depth0pt\special{pdf:bcontent}\special{pdf:code q 0 g 2.453333 4.600000 m 2.146667 4.293333 l 2.453333 3.986667 l 2.453333 3.373333 l 1.840000 2.760000 l 1.226667 2.760000 l 1.226667 2.453333 l 0.613333 1.840000 l 0.306667 1.840000 l 0.000000 1.533333 l 0.000000 0.613333 l 0.613333 0.000000 l 3.373333 0.000000 l 3.986667 0.613333 l 3.986667 0.920000 l 4.293333 1.226667 l 4.293333 1.533333 l 3.986667 1.840000 l 3.986667 2.146667 l 3.680000 2.146667 l 3.373333 2.453333 l 3.373333 3.066667 l 3.680000 3.373333 l 3.680000 3.986667 l 3.066667 4.600000 l h 1.226667 2.146667 m 1.533333 2.453333 l 1.840000 2.453333 l 2.453333 1.840000 l 2.453333 1.226667 l 1.840000 0.613333 l 1.226667 1.226667 l h f* Q}\special{pdf:econtent}\kern4.293333bp}}

\hypertarget{native-566-lower}{}
\begin{center}\small Göttingische gelehrte Anzeigen. Stück 27. Seite 269..272. 14. Februar 1814.\\[3pt]\rule{.72\linewidth}{.3pt}\end{center}

\emph{Philosophical Transactions of the Royal Society of London} for the Jear 1813. IV und 304 Seiten. Meteorological Journal und 8 Seiten Register im Quartformat.

\emph{Beobachtung der Sommersonnenwenden von} 1812 \emph{und} 1813 \emph{und der Wintersonnenwende von} 1812 \emph{auf der Königlichen Sternwarte in Greenwich}, von J. \textsc{Pond} (drei Aufsätze). Das Instrument für diese Beobachtungen ist ein von \textsc{Troughton} gefertigter Mauerkreis von sechs Fuß, den Herr \textsc{Pond} später ausführlich zu beschreiben verspricht. Aus den wenigen Angaben, die hier über ihn mitgeteilt werden, geht hervor, dass seine Besonderheit darin besteht, die Bögen des Himmelsmeridians ohne Bezug zum Zenit zu messen. Deshalb hat er weder ein Bleilot noch eine Libelle. Die Bestimmung der Lage des Zenitpunkts im Meridian bleibt allein dem Zenitsektor vorbehalten. Die hier gerühmte Genauigkeit der Messungen mit diesem Instrument, das einer über Paris erhaltenen privaten Nachricht zufolge 1600 Guineen kosten soll, scheint für diesen Teil der praktischen Astronomie eine neue Epoche zu versprechen. Die Tabelle gibt die Ergebnisse dieser Bestimmungen für die mittlere Schiefe der Ekliptik an.

\begin{center}\fontsize{9}{11}\selectfont\setlength{\tabcolsep}{3pt}\renewcommand{\arraystretch}{1.18}\begin{tabular}{@{}l c r@{\;}r@{\;}r@{}}
Sommersonnenwende von 1812 & $.\;.\;.\;.\;.\;.\;.\;.\;.\;.$ & $23^\circ$ & $27'$ & $50''5$\\
Wintersonnenwende von 1812 & $.\;.\;.\;.\;.\;.\;.\;.\;.\;.$ & $23$ & $27$ & $47.35$\\
Sommersonnenwende von 1813 & $.\;.\;.\;.\;.\;.\;.\;.\;.\;.$ & $23$ & $27$ & $49.5$\\
\end{tabular}\end{center}

{\rightskip=0pt plus 3em Bei der Berechnung der Beobachtungen verwendete Herr \textsc{Pond} die Refraktionstafeln von \textsc{Bradley}.\par}

\emph{Verzeichnis der Nordpolabstände von} 44 \emph{der bedeutendsten Fixsterne} für den Beginn des Jahres 1813 und in einem späteren Aufsatz \emph{ein ähnliches Verzeichnis für} 84 \emph{Fixsterne}, von J. \textsc{Pond}. Das erste dieser \SourceLeaf{567}Verzeichnisse war auch in Deutschland schon früher bekannt geworden. Das zweite ist als erweiterte und verbesserte Ausgabe anzusehen. Die Zahl der Beobachtungen, auf denen diese Bestimmungen beruhen, reicht bei einigen Sternen bis 100, beim Polarstern sogar bis 200. Bei einigen der bedeutendsten Sterne gibt Herr \textsc{Pond} zugleich die Ergebnisse der einzelnen Beobachtungsreihen an (jeweils zehn Beobachtungen). Ihre Übereinstimmung zeugt von der Vortrefflichkeit des Instruments und übertrifft alles bisher in dieser Art Bekannte. Bei vielen Sternen hielt Herr \textsc{Pond} seine Bestimmung bis auf $\frac{1}{4}$ Sekunde für zuverlässig, soweit sie allein vom Instrument abhängt. Eine Auswahl von 30 der bedeutendsten Sterne nördlich des Äquators hat Herr \textsc{Pond} noch in einem eigenen Verzeichnis zusammengestellt, das er Normalkatalog (Standard catalogue) nennt. Als Gegenstück zu \textsc{Maskelyne}s Katalog für die Rektaszensionen will er es immer weiter vervollkommnen.

Mit der praktischen Astronomie hängen auch zwei Aufsätze von \textsc{William Hyde Wollaston} zusammen. Der erste beschreibt ein Verfahren zum Ziehen äußerst feiner Metalldrähte. Dazu bringt man einen Draht aus Gold oder Platin ins Innere eines zunächst dicken Silberdrahts (indem man entweder den Silberdraht hohl bohrt oder den Platindraht in einer Form mit geschmolzenem Silber umgießt), zieht diesen Doppeldraht auf bekannte Weise aus und löst schließlich das Silber mit warmer Salpetersäure wieder ab. Mit Platin gelangen diese Versuche am besten. So erhielt er Fäden von $\frac{1}{1000}$ und $\frac{1}{8000}$ Zoll Dicke, die sich ausgezeichnet zum Einspannen im Brennpunkt astronomischer Fernrohre eigneten. Diese Fäden sind schon viel dünner als die feinsten Spinnenfäden. (Der Rezensent fand durch Messungen, dass die Dicke eines im Brennpunkt eines 16-Zoll-Fernrohrs gespannten Spinnenfadens einen Bogen von $7''8$ abdeckte. Demnach wäre dieser Faden $\frac{1}{1630}$ Zoll dick.) Herr \textsc{Wollaston} ging jedoch noch viel weiter und stellte Fäden von fast unglaublicher Feinheit her (wenn auch nur in kurzen Stücken). Für ihre Dicke waren sie immer noch sehr fest. Ein Platinfaden von $\frac{1}{18000}$ Zoll Dicke trug noch $1\frac{1}{\ObservedWollastonDenominator}$ Gran. Um die Dicke so äußerst feiner Gegenstände mit einer gewissen Genauigkeit messen zu können, bediente er sich einer besoneern Vorrichtung, die er im anderen Aufsatz unter dem Namen Single-lens micrometer beschreibt. Ihr wesentlicher Teil ist eine Linse mit sehr kurzer Brennweite ($\frac{1}{12}$ Zoll), in deren Brennpunkt der Gegenstand gebracht wird. Die Öffnung dieser Linse ist so klein, dass neben der Linse in dem Messingplättchen, in das sie gefasst ist, in einem Abstand von $\frac{1}{23}$ Zoll vom Mittelpunkt ein kleines Loch angebracht werden kann. Die Pupille erhält zugleich Licht vom Gegenstand durch die Linse und von einer verschiebbaren Skala durch das Loch. Die Skala wird so weit entfernt, bis der Gegenstand genau einen oder einige ihrer Teilstriche abdeckt. Der Abstand der Skala und die Zahl der Teilstriche bestimmen dann die Größe des Gegenstands.

Hierher gehört auch \emph{Über die Lichtstärke des Cassegrain-Teleskops im Vergleich mit dem Gregory-Teleskop} von H. \textsc{Kater}. Zwei Paare solcher Teleskope, die derselbe Instrumentenbauer (Herr \textsc{Crickmore} in Ipswich) angefertigt hatte, wurden miteinander verglichen. Dabei wurden der Unterschied im Flächeninhalt der großen Spiegel und der Verlust durch die kleineren Spiegel samt ihren Tragarmen berücksichtigt. Bei dem einen Paar verhielt sich die Lichtstärke wie 7 zu 3, beim anderen wie 3 zu 2, beide Male zugunsten des \textsc{Cassegrain}-Teleskops. Beim ersten Paar hatte der Spiegel des \textsc{Cassegrain}-Teleskops den Vorteil einer etwas besseren Politur, beim anderen der des \textsc{Gregory}-Teleskops. Daher glaubt der Verfasser, dass das \textsc{Cassegrain}-Teleskop unter ganz gleichen Umständen etwa doppelt so viel Licht hat wie das \textsc{Gregory}-Teleskop. Wenn wir voraussetzen dürfen, dass die Öffnung, an die das Auge unmittelbar gehalten wurde, in allen vier Teleskopen groß genug war, um das gesamte Licht der großen Spiegel durchzulassen, — der Verfasser geht auf diesen Punkt nicht ein, doch der Rezensent \SourceLeaf{568}hat mehrere \textsc{Gregory}-Teleskope in Händen gehabt, bei denen diese Öffnung zu klein war — dann sind diese Erfahrungen sowohl für den Astronomen als auch für den Physiker höchst bemerkenswert. Man wäre geneigt, der Erklärung des Verfassers zuzustimmen. Er sucht den Grund der sonderbaren Erscheinung darin, dass im \textsc{Gregory}-Teleskop zwei reelle Bilder entstehen, im \textsc{Cassegrain}-Teleskop dagegen nur eines. Daraus schließt er, dass die Kreuzung der Strahlen in einem Punkt oder die Entstehung eines reellen Bildes einen gewissen Lichtverlust bewirkt. Er glaubt, damit hänge vielleicht die Erfahrung zusammen, dass \textsc{Galilei}-Fernrohre unter gleichen Umständen mehr Wirkung erzielen als astronomische Fernrohre.
